\section{一次改动怎样通过测试和审查}

以“给 Stage 1 的 ATA 等待增加超时”为例。纯函数测试可以检查预算递减和状态
分类，镜像测试可以检查加载范围，QEMU 故障场景可以让 BSY 始终不清零，正常
启动则确认真实镜像没有被过短预算误伤。四种检查针对不同错误，不必强行合并
为一个庞大的整机测试。

\topic{测试从便宜而精确的层次开始}
\begin{osgridlongtable}{L{0.20\textwidth}L{0.32\textwidth}L{0.36\textwidth}}
\textbf{层次} & \textbf{适合发现} & \textbf{例子}\\ \hline
单元测试 & 纯状态转换、边界和错误枚举 & ELF 截断、地址溢出、引用计数\\ \hline
随机/性质测试 & 大量操作组合后的不变量 & 分配释放、VMA 区间、journal 模型\\ \hline
产物审计 & 链接和打包后的字节事实 & 复位向量、ROM 尺寸、ELF 段与磁盘 LBA\\ \hline
QEMU 集成 & CPU 模式、端口、中断和真实控制流 & 长模式、IRQ、Ring 3、ATA\\ \hline
整机场景 & 多模块组合与用户可见行为 & fork/exec、流水线、重启后文件读取\\ \hline
\end{osgridlongtable}

\topic{故障注入要检查失败后的状态}
让第四次页帧申请失败，断言前三帧已经释放，页表没有留下半条映射，统计也回到
调用前。让 journal 在 commit 前中断，重启后应丢弃事务；在有效 commit 后
中断，则应重放。只断言“函数返回失败”会漏掉最危险的资源泄漏与半状态。

\topic{随机测试必须能够复现}
测试失败时打印 seed、操作序号和缩减后的输入。磁盘崩溃测试还要记录中断点、
镜像摘要和第一次不一致的块。没有这些信息，下一次运行不再失败时就失去了
定位机会。

\topic{整机测试必须有超时和最后输出}
QEMU 可能因为三重故障重启，也可能在设备轮询、锁或等待队列中永久停住。
驱动脚本为每个场景设置上限，超时后保存串口尾部、QEMU 退出状态和必要的
诊断文件。测试进程也要被清理，避免一个失败场景占着磁盘或端口影响后续测试。

\topic{技术审查沿四类事实展开}
\begin{enumerate}[label=\arabic*.]
  \item 对照当前源码，确认名称、顺序、容量和日志仍然存在；
  \item 对照 Intel 手册或设备规范，确认寄存器位、保留位和访问时序；
  \item 对照最终 ROM、ELF 和磁盘镜像，确认链接与打包没有改变事实；
  \item 在 QEMU 中运行成功与故障场景，确认观察结果能够区分各条路径。
\end{enumerate}

阅读时也可以用这四类事实互相校验。若正文、图、运行结果和处理器手册给出
不同结论，先回到具体地址、寄存器位和代码路径逐项核对，不能靠换一种措辞
掩盖差异。找不到来源的精确数字，应回到代码或规格重新计算。

\topic{评审时问具体问题}
看到加法，检查溢出与半开区间；看到裸指针，检查地址空间和生命周期；看到
引用，检查最后释放者；看到等待，检查唤醒是否可能丢失以及谁能超时；看到
硬件状态位，检查轮询上限和错误优先级；看到日志，检查它是在状态改变之前还是
之后打印。这样评审会落在真实行为上，而不会停在“结构很完整”的评价。

\topic{完成一项修改后的学习记录}
最后用一页纸回答：运行前知道什么，哪条现象促使修改，第一次猜测是什么，
哪项检查证实或否定了它，代码改了哪条状态转换，失败路径怎样验证。真实的
思考过程能帮助后来者复现，也会暴露作者尚未理解的空白。
